% Created 2026-09-05 Sat 13:25
% Intended LaTeX compiler: xelatex
\documentclass[12pt]{article}
\usepackage[utf8]{inputenc}
\usepackage[T1]{fontenc}
\usepackage{graphicx}
\usepackage{longtable}
\usepackage{wrapfig}
\usepackage[margin=1in]{geometry}
\usepackage{rotating}
\usepackage{setspace}
\onehalfspacing
\usepackage[normalem]{ulem}
\usepackage{amsmath}
\usepackage{parskip}
\usepackage{amssymb}
\usepackage{capt-of}
\setcounter{secnumdepth}{4} % Tells LaTeX to number down to paragraphs
\setcounter{tocdepth}{4}
\usepackage{hyperref}
\usepackage[style=oscola, backend=biber]{biblatex}
\addbibresource{/Users/krishnajani/Documents/Libib.bib}
\usepackage{csquotes}
\usepackage[british]{babel}
\setlength{\parskip}{0.5\baselineskip}
\author{Krishna Jani}
\date{\today}
\title{CoC of Essar Steel v. Satish Kumar Gupta}

\begin{document}

\maketitle

\section*{Facts}

\begin{itemize}
	\item CIRP proceedings were instituted by State Bank of India and Standard Chartered Bank against SR Steel, in respect of which one Satish Kumar Gupta was appointed as Resolution Professional (RP).
	\item Two resolution plans were submitted, one by Arcelor Mittal Private Limited and the other by NuMetal Private Limited, both of which were found to be ineligible under Section 29A of the IBC. Since both Arcelor Mittal and Numetal had not paid off their NPAs to their respective corporate debtors.
	\item However, an order was passed under Article 142 of the Constitution directing both parties to pay off their NPAs and resubmit the resolution plan.
	\item Arcelor Mittal resubmitted the resolution plan, which was evaluated by the COC on the same date, and the resolution plan was thereafter approved by a 92\% majority of the CoC.
	\item Interestingly, the NCLT disposed of the application to allow the resolution plan on the grounds that it did not align with the principles of equity and fair play, since the dues of the operational creditors did not get similar treatment as compared to the dues of the financial creditors.
	\item The NCLT recommended that 85\% of the amount offered by the SRA should be distributed among the FCs, while the remaining 15\% should be distributed among the rest of the OCs. Doing so would ensure that at least 50\% of the admitted claims of the OCs can be paid off. If the FCs are willing to sacrifice the interest component on their principal, then the debts of the OCs could be satisfied in a reasonable and fair manner.
	\item The NCLAT directed the COC to take a decision on the suggestions that were made by the NCLAT. The COC decided by a majority of 70.73\% to make an ex-gratia payment of 1,000 crores to operational creditors above INR 1 Cr.
	\item By its final judgement, the NCLAT held:
		\begin{itemize}
			\item There can be no difference between financial creditor and operational creditor in the payment of dues, and therefore FC and OC must deserve equal treatment under the resolution plan. Therefore, the proceeds payable are redistributed to ensure that FCs and OCs are paid 67.60\% of their admitted claims.
			\item Security interest is irrelevant at the stage of resolution for the purpose of allocation of payments.
			\item If each FC with a claim more than INR 10 lakhs is to be paid 60.7\% of its admitted claim, 3 OCs, by definition, have a separate class within themselves and can be subdivided for the purpose of distribution. Therefore, ACs/OCs with a claim greater than or equal to INR 1 Cr shall be paid 60.2\% of their admitted claims.
			\item Profits for profits generated by CD during the CIRP shall be equally distributed among FCs and OCs.
			\item A subcommittee or core committee cannot be constituted under the IBC, and the COC must take all its decisions themselves.
		\end{itemize}
\end{itemize}

\section*{Appellant's Arguments}
\begin{itemize}
	\item The Code does not mandate identical treatment of differently situated creditors, either inter se within financial creditors or financial creditors vis-à-vis operational creditors. The Code only posits equal treatment of different classes of creditors that are similarly placed. FCs, as a class, have a superior status as against OCs. The same is the case with secured creditors vis-à-vis unsecured creditors.
	\item The classification of creditors based on the nature of the debt or security interest is \textit{sine qua non} to any insolvency code. If secured financial creditors are to be treated at par with unsecured creditors, such secured creditors would rather vote for liquidation than corporate resolution, contrary to the main objective that is sought by the Code. Since secured creditors are at the top of the waterfall mechanism and will often recoup more of their investment as against CIRP.
	\item Therefore, while other secured creditors were secured to the extent of 99.66\% of their outstanding dues, Standard Chartered Bank was secured merely by a pledge of shares held by the CD in an offshore Mauritian subsidiary. The fair value of which was determined at mere 0.7\% of the admitted claim of Standard Chartered Bank. Therefore, Standard Chartered was an unsecured creditor. The differential treatment of which is not wholly unjustified, since there exists a clear differentia between secured and unsecured creditors within the IBC CIRP mechanism.
	\item There is a judicial hands-off approach when it comes to the commercial wisdom of the COC, which has been directly infringed by the impugned judgement. A resolution plan is a consent-based plan proposed by the SRA for the CD. The counterparty to such a plan is the COC, which is required to give a minimum consent of 66\% voting share, which is the basis for the NCLT to approve a resolution plan. Any modification done by the NCLT of such a plan is illegal.
	\item The COC has the jurisdiction to deal with all commercial aspects of a resolution plan, and in this regard, setting up of a subcommittee or a core committee within the COC is permissible under the Code. Since no decision-making power is delegated to such a committee, nor does it decide or recommend the distribution of amounts

\end{itemize}

\section*{Respondent's Arguments}

\begin{itemize}
	\item Arcelor Mittal made an offer of INR 42,000 crore before the Supreme Court. Yet, what ultimately turned out was a payment of a lesser value: INR 39,500 as upfront, and INR 2,500 crore was added towards guaranteed working capital adjustment.

	\item It was submitted that this adjustment is, in fact, a result of backdoor negotiation between the majority in the core committee of the COC and the SRA, where CoC agreed to reduce the upfront offering if it would acquire the debts of OSPIL (which was a wholly owned subsidiary owning a slurry pipeline in Odisha).
	\item It was contended that the formation of a core committee/subcommittee was against the provisions of the Code, since it went on to conduct various secret negotiations with the SRA, which resulted in the burial of Standard Chartered Bank's debt almost entirely. Its entitlement was reduced from 2,585 crore to the value of the underlying security, which is, as previously stated, merely 0.7\% of the advertised claim.
	\item The resolution plan of Arcelor Mittal was also repugnant to Regulation 38(1)(a), since it did not deal with the interests of all stakeholders. The role of the COC within the Code is limited to considering the feasibility and viability of the plan and does not include the manner of distribution of the amount payable to the resolution applicant. Therefore, the decision of the COC with respect to the matter of distribution and disbursement of the said sums is illegal and arbitrary.
	\item Once the creditor is classified as an FC, he is entitled to equal treatment with all other FCs, irrespective of whether it is secured or unsecured. The word "secured creditor", as defined by Section 3(30) of the Code, is used first in Section 53, which is contained in the chapter titled "Liquidation Process", and not in Chapter 2, which is titled "Corporate Insolvency Resolution Process". Therefore, the differential treatment between secured and unsecured creditors is permissible only during liquidation proceedings and not during CIRP proceedings. No categorization can be based on the security interest of the financial creditor.
\end{itemize}

\section*{Court's Observations}

\begin{itemize}
	\item The Court begins with a discussion of the role of the RP and the SRA to lay ground for CoC. For RP's the Court observes that it aids in
		\begin{itemize}
			\item Managing the affairs of the CD as a going concern during the CIRP
			\item Appoint and convene the meetings of the CoC
			\item Create a detailed memo of the CD and providing for EOI's to facilitate IRP
			\item Collect, collate and admit the claims of creditors.
			\item Examine the RP to ensure that it is complete in all aspects before it is placed before the CoC
		\end{itemize}
		For PRA (Prospective Resolution Applicant) it is stated that
		\begin{itemize}
			\item The RP submitted by the PRA must provide for the measures necessary for the insolvency resolution of the CD and for the maximisation of the value of its assets.
			\item It is allowed to provide for either the satisfaction or modification of the security interest of the secured creditor and also may provide for a modification in the amount payable to various class of creditors.
		\end{itemize}
	\item 
\end{itemize}
\end{document}
